\documentclass[12pt]{article}
\usepackage{fullpage}
\usepackage{amsmath,amssymb,amsthm}

\newtheorem{theorem}{Theorem}
\newtheorem{lemma}{Lemma}

\begin{document}

\title{A weighted Bernoulli mean-tail inequality}
\author{}
\date{}
\maketitle

\section*{Problem statement and interpretation}

I interpret the question in the universal sense: determine those \(p\in[0,1]\)
for which
\[
   \Pr\!\left[\sum_{i=1}^m w_i v_i\ge p\right]\ge p
\]
holds for every \(m\ge1\) and every choice of non-negative weights
\(w_1,\ldots,w_m\) with \(\sum_i w_i=1\), where the \(v_i\)'s are i.i.d.\
\({\rm Bernoulli}(p)\).

\section*{The answer}

The required values are
\[
        \boxed{\ [0,1/3]\ \cup\ \{1/2,1\}\ } .
\]

We shall use the following finite form of a theorem of Kellerer.  It is
often called Kellerer's weighted Bernoulli inequality.

\begin{lemma}[Kellerer]\label{lem:kellerer}
Let \(0\le p\le 1/3\).  If \(w_1,\ldots,w_m\ge0\) and
\(\sum_i w_i=1\), and if \(v_1,\ldots,v_m\) are i.i.d.\
\({\rm Bernoulli}(p)\), then
\[
   \Pr\!\left[\sum_{i=1}^m w_i v_i\ge p\right]\ge p .
\]
\end{lemma}

For reference, this is the specialization to finite weighted sums of
Kellerer's ``analogue domains'' theorem: the lower \(p\)-biased measure
of a weighted halfspace whose threshold is its \(p\)-mean is \(p\) for
\(p\le1/3\).  The statement is exact for arbitrary real weights; no
limiting or equal-weight approximation is involved.

We now prove the characterization.

\subsection*{Sufficiency}

If \(0\le p\le1/3\), the desired inequality is precisely
Lemma~\ref{lem:kellerer}.  The endpoint \(p=0\) is also immediate, since
the event is then the whole sample space.

Next let \(p=1/2\).  Put
\[
       X=\sum_{i=1}^m w_i v_i .
\]
Since \(1-v_i\) has the same distribution as \(v_i\), the random
variables \(X\) and \(1-X\) have the same distribution.  Hence
\[
 \Pr[X<1/2]=\Pr[1-X<1/2]=\Pr[X>1/2].
\]
Writing \(a=\Pr[X>1/2]\) and \(b=\Pr[X=1/2]\), we get
\(2a+b=1\), and therefore
\[
      \Pr[X\ge1/2]=a+b=\frac{1+b}{2}\ge\frac12 .
\]
Finally, for \(p=1\), all \(v_i=1\) almost surely, so
\(\sum_i w_i v_i=1\) almost surely and the probability is \(1\).

Thus every \(p\in[0,1/3]\cup\{1/2,1\}\) works.

\subsection*{Necessity}

It remains to show that no other \(p\) can work.

First suppose
\[
        \frac13<p<\frac12 .
\]
Take \(m=3\) and \(w_1=w_2=w_3=1/3\).  Then
\[
 \sum_{i=1}^3 w_i v_i=\frac{v_1+v_2+v_3}{3}.
\]
Because \(p>1/3\), the event
\(\frac{v_1+v_2+v_3}{3}\ge p\) requires at least two successes.  Hence
\[
\Pr\!\left[\frac{v_1+v_2+v_3}{3}\ge p\right]
  =\Pr[{\rm Bin}(3,p)\ge2]
  =3p^2(1-p)+p^3
  =3p^2-2p^3 .
\]
But
\[
  3p^2-2p^3-p
   =p(1-p)(2p-1)<0
\]
for \(1/3<p<1/2\).  Thus the required inequality fails throughout
\((1/3,1/2)\).

Now suppose
\[
        \frac12<p<1 .
\]
Take \(m=2\) and \(w_1=w_2=1/2\).  Then
\[
 \sum_{i=1}^2 w_i v_i=\frac{v_1+v_2}{2}.
\]
Since \(p>1/2\), the event
\(\frac{v_1+v_2}{2}\ge p\) requires two successes.  Therefore
\[
\Pr\!\left[\frac{v_1+v_2}{2}\ge p\right]=p^2<p,
\]
so the inequality fails for every \(p\in(1/2,1)\).

Combining the sufficiency and the two counterexample ranges gives exactly
\([0,1/3]\cup\{1/2,1\}\).

\begin{thebibliography}{1}

\bibitem{Kellerer1963}
H.~G. Kellerer,
\emph{Zur Existenz analoger Bereiche},
Zeitschrift f\"ur Wahrscheinlichkeitstheorie und Verwandte Gebiete
\textbf{1} (1963), 240--246.

\end{thebibliography}

\end{document}
